% Software Requirements Specification — Vehicle Dynamics (IEEEtran)
% Requires IEEEtran.cls (TeX Live / MiKTeX IEEE package).

\documentclass[conference]{IEEEtran}
\usepackage{amsmath,amssymb,amsfonts}
\usepackage{booktabs}
\usepackage{cite}
\usepackage{url}
\usepackage{hyperref}

\hyphenation{op-tical net-works semi-conduc-tor}

\begin{document}

\title{Software Requirements Specification\\for a Vehicle Dynamics Simulator with DEM Soft Soil}

\author{\IEEEauthorblockN{Sankalp Amaravadi}
\IEEEauthorblockA{Vehicle Dynamics Project\\
Email: (TBD)}}

\maketitle

\begin{abstract}
This Software Requirements Specification (SRS) defines functional and non-functional
requirements for a greenfield vehicle-dynamics backend. The system couples a
multi-body vehicle model to discrete-element-method (DEM) loose-ground simulation
with a target scale of millions of particles on GPU. A Python orchestration layer
provides configuration, scenario control, and data export. Graphical front-end
requirements are deferred.
\end{abstract}

\begin{IEEEkeywords}
vehicle dynamics, DEM, soft soil, SRS, GPU simulation, Python
\end{IEEEkeywords}

\section{Introduction}

\subsection{Purpose}
This document specifies software requirements for the Vehicle Dynamics simulator
backend. It is intended for developers, reviewers, and stakeholders who validate
scope before design and implementation proceed.

\subsection{Document Conventions}
Requirements are labeled FR-$n$ (functional) and NFR-$n$ (non-functional). Items
marked TBD remain open until later phases.

\subsection{Scope}
The product is a command-line / API-driven simulator that:
\begin{itemize}
  \item advances vehicle dynamics on terrain;
  \item represents loose ground as DEM particle assemblies;
  \item couples tire--soil interaction;
  \item exports time-series and logs for analysis.
\end{itemize}
User-interface clients are out of scope for early phases.

\subsection{Definitions and Acronyms}
\begin{itemize}
  \item \textbf{DEM} --- Discrete Element Method for granular media.
  \item \textbf{Soft soil / loose ground} --- Terrain modeled as many DEM bodies.
  \item \textbf{SRS} --- Software Requirements Specification.
  \item \textbf{SDD} --- Software Design Description.
\end{itemize}

\section{Overall Description}

\subsection{Product Perspective}
The backend is greenfield software. It is not derived from the CS4632 dealership
discrete-event simulation codebase. Long-term architecture separates Python
orchestration from GPU DEM kernels.

\subsection{Product Functions}
Primary capabilities include vehicle motion simulation, large-scale DEM soil,
tire--soil coupling, scenario-driven runs, and persistent outputs for
post-processing.

\subsection{User Characteristics}
Primary users are developers and analysts executing batch or interactive runs
from the command line. Secondary users may consume exported data in external
tools.

\subsection{Constraints}
\begin{itemize}
  \item Orchestration language: Python~3.
  \item Million-scale DEM requires GPU-capable hardware and drivers.
  \item Documentation artifacts (SRS, SDD, testing) reside under
        \texttt{documentation/}.
\end{itemize}

\section{Functional Requirements}
\begin{itemize}
  \item \textbf{FR-1:} The system shall provide a runnable entry point
        (\texttt{main.py}).
  \item \textbf{FR-2:} The system shall accept configuration for vehicle and
        soil parameters (TBD schema).
  \item \textbf{FR-3:} The system shall advance a coupled vehicle--DEM
        simulation over a configured time horizon (TBD).
  \item \textbf{FR-4:} The system shall persist simulation outputs suitable for
        post-processing (TBD formats).
  \item \textbf{FR-5:} The system shall support scenario definitions for
        repeatable maneuvers (TBD library).
\end{itemize}

\section{Non-Functional Requirements}
\begin{itemize}
  \item \textbf{NFR-1:} The DEM path shall target GPU execution for
        million-scale particle counts.
  \item \textbf{NFR-2:} Hot paths may use GPU kernels while remaining
        invocable from Python orchestration.
  \item \textbf{NFR-3:} Source and documents shall remain suitable for version
        control and reproducible builds.
  \item \textbf{NFR-4:} Small-scale DEM modes shall exist for development and
        continuous integration without requiring full GPU scale.
\end{itemize}

\section{Requirements Reserved for Later Phases}
Detailed vehicle degrees of freedom, tire constitutive model, DEM contact law,
validation metrics, and scenario catalog remain TBD and will be amended in
subsequent SRS revisions.

\section*{Acknowledgment}
This document follows the IEEE conference \LaTeX\ template (\texttt{IEEEtran}).

\begin{thebibliography}{00}
\bibitem{ieee830} IEEE Std 830-1998 (withdrawn), \emph{IEEE Recommended Practice
for Software Requirements Specifications}.
\bibitem{ieeetran} M.~Shell, ``IEEEtran homepage,''
\url{https://www.michaelshell.org/tex/ieeetran/}.
\end{thebibliography}

\end{document}